% 张三慧V2第24章16题独立解析候选。主教材PDF313--315，配套S328--333。
% 原图精确重绘与独立学术审查仍待办，保留源书起偏角/临界角与偏振轴差异。
\section*{第24章\quad 光的偏振：独立解析}
\paragraph{习题 24.1}
设入射自然光强度 \(I_0\)。
通过第一偏振片变为 \(I_0/2\)，
原第二片与第一片成60°，故
\(I_1=(I_0/2)\cos^260^\circ=I_0/8\)。
加入中间片使两相邻夹角各30°，
\[
I'=\frac{I_0}{2}\cos^230^\circ\cos^230^\circ
=\frac{9I_0}{32}
=\boxed{\frac94I_1=2.25I_1}.
\]
插入第三片反而能提高通过交叉较大夹角偏振片的能流，
因为中间投影改变了最终可透分量。

\paragraph{习题 24.2}
马吕斯定律 \(I=(I_0/2)\cos^2\theta\)。
最大透射光强 \(I_{\max}=I_0/2\)。
(1) \(I=I_{\max}/3\)：
\[
\boxed{\theta=\arccos\sqrt{1/3}\simeq54.74^\circ}.
\]
(2) \(I=I_0/3\)：
\[
\boxed{\theta=\arccos\sqrt{2/3}\simeq35.26^\circ}.
\]
不可混淆“最大透射强度”与“入射光强度”。

\paragraph{习题 24.3}
只看与传播方向垂直的平面，
第二片旋转后的\emph{实际有效透振轴}
同第一片的夹角记作 \(\alpha\)。
由原图24.38对第二偏振轴先绕光轴转
\(\theta\)，再绕竖直轴转 \(\varphi\) 的投影几何，
\[
\cos\alpha=
\frac{\cos\theta}
{\sqrt{\cos^2\theta+\sin^2\theta\cos^2\varphi}}.
\]
因此，排除投影轴退化的极限几何时
\[
\boxed{I_2=\frac{I_0}{2}\cos^2\alpha
=\frac{I_0\cos^2\theta}
{2(\cos^2\theta+\sin^2\theta\cos^2\varphi)}}.
\]
\(\varphi=0\) 时退化为标准 \(I_0\cos^2\theta/2\)；
直接将 \(\cos^2\theta\cos^2\varphi\) 相乘
通常不满足横向投影物理。

\paragraph{习题 24.4}
点号表示电矢量垂直于入射面（s偏振），
短线表示在入射面内（p偏振）。
同一各向同性平面界面可分别对两线偏振分量
应用菲涅耳反射、透射系数；
反射角等于入射角，折射角满足
\(n_1\sin i=n_2\sin t\)。
对于非临界全反射情形：
\[
r_s=\frac{n_1\cos i-n_2\cos t}
{n_1\cos i+n_2\cos t},\quad
r_p=\frac{n_2\cos i-n_1\cos t}
{n_2\cos i+n_1\cos t}.
\]
在 \(i=i_b\) 时 \(\tan i_b=n_2/n_1\)，
于是 \(\boxed{r_p=0}\)，但一般 \(r_s\ne0\)。
因此题图(a)--(e)按入射标志逐个绘图时：
纯s入射的反射/折射光仍为s；
纯p入射在 \(i\ne i_b\) 时反射和透射仍为p，
在 \(i_b\) 时无p反射；
混合入射且 \(i=i_b\) 时反射只含s，
透射同时含s与p，一般为部分线偏振态。
原图每一箭头的具体点/线布局及反射相位
尚需逐幅TikZ复核，不能仅据文字宣称图形认证。

\paragraph{习题 24.5}
水 \(n_1=1.33\) 到玻璃 \(n_2=1.50\)：
\[
\boxed{i_{B,w\to g}
=\arctan\frac{1.50}{1.33}\simeq48.43^\circ}.
\]
玻璃到水：
\[
\boxed{i_{B,g\to w}
=\arctan\frac{1.33}{1.50}\simeq41.57^\circ}.
\]
由于两折射率比互为倒数，
\(\boxed{i_{B,w\to g}+i_{B,g\to w}=90^\circ}\)。
布儒斯特角处的反射光只有s偏振分量。

\paragraph{习题 24.6}
题给的是\emph{全反射临界角} \(\beta=45^\circ\)，
而不是布儒斯特角。
由高折射率一侧射向低折射率一侧，
\[
\sin\beta=\frac{n_2}{n_1}=\sin45^\circ.
\]
从\emph{同一高折射率侧}入射的起偏角满足
\(\tan i_B=n_2/n_1\)，因此
\[
\boxed{i_B=\arctan(\sin45^\circ)
\simeq35.26^\circ}.
\]
不能直接用45°当作起偏振角。

\paragraph{习题 24.7}
由空气到材料的布儒斯特关系
\(\tan i_B=n/n_{\rm air}\)，取 \(n_{\rm air}\simeq1\)，
\[
\boxed{n\simeq\tan58^\circ\simeq1.60}.
\]

\paragraph{习题 24.8}
湖面将入射太阳光反射为纯s线偏振光
需在布儒斯特角入射：
\(i_B=\arctan(1.33/1)\simeq53.06^\circ\)。
太阳的仰角（与水平面的夹角）为
\[
\boxed{\alpha=90^\circ-i_B\simeq36.94^\circ}.
\]
这是镜面反射和理想纯净水面的几何近似。

\paragraph{习题 24.9}
棱镜正三角形截面的斜面与底边夹60°，
欲使 \(e\) 光棱镜内沿底边传播，
则它相对入射面法线的折射角
\(r_e=90^\circ-60^\circ=30^\circ\)。
光轴垂直图面，该方向的有效折射率
取源书 \(n_e=1.4864\)，于是
\[
\boxed{i=\arcsin(n_e\sin30^\circ)
=\arcsin(1.4864/2)\simeq48.0^\circ}.
\]
另一束 \(o\) 光折射角
\[
r_o=\arcsin\frac{\sin i}{n_o}
=\arcsin\frac{\sin48^\circ}{1.6584}
\simeq26.6^\circ.
\]
因此在棱镜内比平行底面的 \(e\) 光
稍向底边倾斜；到第二界面时需再用
\(n_o\sin i'_o=\sin t'_o\)
求出射方向（源书给约 \(65.9^\circ\)）。
由于光轴指向图外，
\(o,e\) 分别为互相正交的线偏振光，
具体点/线箭头需依原图24.40作图验证。

\paragraph{习题 24.10}
原图24.41的自然光入射后，先在光轴沿
\(AB\) 的第1个方解石棱镜内
分解为相互正交的普通光 \(o\) 和非常光 \(e\)。
两束光到达 \(BD\) 斜界面时，因为第2块
棱镜光轴改成垂直图面，
同一振动方向对于另一个晶体的
\(o/e\) 身份交换；波法线满足各自相应的
折射定律，出现两束分离的出射光。
图示中一束向右上偏、一束向右下偏；
两束的线偏振方向正交。
简单用相同的标量折射率跨过内部界面，
会漏掉双折射与光轴交换。
\textit{精确出射角和各箭头仍须根据图24.41/配套S PDF330--331重画审核。}

\paragraph{习题 24.11}
题目中的 \(P_1,P_2\) \emph{互相正交}，
不是平行偏振片。第一片透光：
\[
\boxed{I_1=I_i/2,\quad\text{线偏振沿 }P_1}.
\]
四分之一波片不吸收，总强度不变，
但两晶体本征轴方向分量
\((A\cos60^\circ,A\sin60^\circ)\)
有相对相位差 \(\pi/2\)：
\[
\boxed{I_2=I_i/2,\quad\text{椭圆偏振}}.
\]
最后一片 \(P_2\) 与 \(P_1\) 正交；
透光沿垂直于 \(P_1\) 的方向，单一相干电场
分量投影强度
\(I_3=I_1\frac12\sin^2(2\cdot60^\circ)\)，因此
\[
\boxed{I_3=\frac{3}{16}I_i,\quad
\text{线偏振沿 }P_2}.
\]
若错误把 \(P_2\) 设为平行，
会得到 \(5I_i/16\)，与原图不同。

\paragraph{习题 24.12}
相位延迟条件
\(|n_o-n_e|\,h_{\min}=\lambda/4\)：
\[
\boxed{h_{\min}=\frac{632.8\,\mathrm{nm}}
{4(1.66-1.49)}
\simeq931\,\mathrm{nm}}.
\]
此为沿波片法线正入射的最小厚度；
为有两个不同横向本征轴，
晶体光轴必须取\emph{平行于波片表面}
（并与光传播方向垂直）的方向。

\paragraph{习题 24.13}
忽略色散，原片对800nm的延迟光程
\(\Delta=(n_o-n_e)h=800/4=200\,\mathrm{nm}\)。
对400nm入射，其相位延迟
\[
\delta=\frac{2\pi\Delta}{400\,\mathrm{nm}}=\pi.
\]
故对400nm光此片是\emph{半波片}而不是四分之一波片。
入射线偏振与本征轴夹45°，
两正交等幅分量经相位差 \(\pi\) 合成仍线偏振，
但振动方向关于本征轴镜像，
\[
\boxed{\text{出射仍为线偏振，
且振动方向与原线偏振正交}}.
\]

\paragraph{习题 24.14}
按原图24.43，两块方解石的入射外侧界面
光线相对法线为 \(22^\circ\)。
对普通光用 \(n_o\simeq1.658\)：
\[
r_o=\arcsin\frac{\sin22^\circ}{1.658}
\simeq13.1^\circ.
\]
几何给普通光在加拿大胶层的入射角
\(i'_o\simeq90^\circ-r_o\simeq76.9^\circ\)。
普通光从方解石到胶层 \(n_g=1.55\)
的全反射临界角
\[
\boxed{\beta_o=\arcsin\frac{1.55}{1.658}
\simeq69.2^\circ<i'_o},
\]
故 \(o\) 光全反射后被涂黑底面吸收。
对于 \(e\) 光，沿题图的有效折射率约
\(n_e=1.486<1.55\)，故跨越胶层
\emph{不会发生从高到低折射率的全反射}，
继续从第2块晶体出射。
\[
\boxed{\text{透过尼科耳棱镜的光为纯 }e
\text{ 线偏振光}}.
\]
本解中的临界角采用源书近似标量折射率；
真实晶体的走离角与光轴方向需进一步图示复核。

\paragraph{习题 24.15}
线偏振光由两等幅相反旋向圆偏振分量合成。
穿过长度 \(L\) 相对相位差
\(\Delta\phi=2\pi(n_L-n_R)L/\lambda\)，
合成线偏振方位角旋转其一半。
故旋光率
\[
\boxed{\left|\frac{d\alpha}{dL}\right|
=\frac{\pi|n_L-n_R|}{\lambda}
=\frac{\pi(1.55821-1.55810)}{396.8\times10^{-9}}
\simeq871\,\mathrm{rad/m}
\simeq49.9^\circ/\mathrm{mm}}.
\]
旋转正负号需明确是沿传播方向观看还是迎光观看；
仅由给定的 \(n_L>n_R\) 不能不说明约定而写绝对左右。

\paragraph{习题 24.16}
选取两等幅反向传播的正交场
\[
\boldsymbol E(x,t)=E_0\cos(\omega t-kx)\,\hat{\boldsymbol y}
+E_0\cos(\omega t+kx)\,\hat{\boldsymbol z}.
\]
相对相位 \(\delta=2kx=4\pi x/\lambda\)。
在 \(x=0\)，同相等幅合成与 \(y\) 轴45°线偏振。
向 \(+x\) 每走 \(\lambda/8\)，相位差增 \(\pi/2\)：
\[
\begin{array}{c|c|l}
x&\delta&\text{偏振态}\\\hline
0&0&+45^\circ\text{线偏振}\\
\lambda/8&\pi/2&\text{圆偏振，第一旋向}\\
\lambda/4&\pi&-45^\circ\text{线偏振}\\
3\lambda/8&3\pi/2&\text{圆偏振，相反旋向}\\
\lambda/2&2\pi&+45^\circ\text{线偏振恢复}
\end{array}
\]
以后沿 \(x\) 以 \(\lambda/2\) 为周期重复。
圆偏振手性取决于沿哪一传播方向观察：
源书S PDF333把沿 \(-x\) 观察
\(\lambda/8\) 定为右旋、\(3\lambda/8\) 定为左旋。
此物理图像是两列反向行波的空间偏振梯度，
不应误作一列前进圆偏振波。
